\readerheading{Why the threshold is one.}
On a gap-free edge, the open traces overlap, so $B_i+A_{i+1}>1$.
If every edge is gap-free, summing gives
\[
 6<\sum_{i=0}^5(B_i+A_{i+1})=\sum_{i=0}^5(A_i+B_i).
\]
At least one V triangle must therefore be supercritical. A nonempty gap
breaks this cyclic argument: it is precisely where the C triangle takes
over a boundary obligation. Equality $B_i+A_{i+1}=1$ still leaves the
common endpoint uncovered by the two open V triangles; it belongs to the
gap branch, not to the overlap branch.

The distinction between actual and selected reaches is also geometric.
The numbers $A_i,B_i,C_i$ describe a fixed triangle. The later lower bounds
$a_i,b_i,c_i$ prescribe points it must contain and may be reduced when making
a common comparison. Reducing a demand does not change the original
triangle or the count $N_+$.
